\subsection{Aggregation \& Grouping}

\subsubsection{What aggregation means (rows $\rightarrow$ groups)}
\J Aggregation summarizes many rows into fewer rows (often one row per group). You typically 
use aggregate functions like \texttt{COUNT}, \texttt{SUM}, \texttt{AVG}, \texttt{MIN}, 
and \texttt{MAX}.

\M Conceptually, aggregation works like this:
\begin{enumerate}
  \item Build the row set (FROM/JOIN).
  \item Filter rows (\texttt{WHERE}).
  \item Group rows (\texttt{GROUP BY}).
  \item Compute aggregates per group.
\end{enumerate}
If you don't use \texttt{GROUP BY}, the whole result is treated as a single group.

\Sx Always be explicit about what your groups represent (business meaning). Before writing 
the query, state it in plain language: ''One row per customer'', ''One row per month'', etc. 
This prevents subtle correctness bugs.

\E Optimizers can compute aggregates using different strategies (hash aggregate vs sort-based 
aggregate). You don't need internals yet, but you should learn to recognize when grouping is 
the expensive part of a query.

\Pitfall
\begin{itemize}
  \item Forgetting what ''one row per group'' means and accidentally returning duplicate groups due 
  to join multiplicity upstream.
\end{itemize}

\subsubsection{GROUP BY correctness rules (and why SQL complains)}
\J When you use \texttt{GROUP BY}, every selected column must be either:
\begin{itemize}
  \item included in the \texttt{GROUP BY}, or
  \item wrapped in an aggregate function.
\end{itemize}

\M This rule exists because, within a group, a non-aggregated column could have multiple values. 
SQL forces you to be explicit so results are deterministic.

\Sx If you find yourself wanting to select a non-grouped column, you usually need one of 
these patterns:
\begin{itemize}
  \item Add the column to \texttt{GROUP BY} (if it truly defines the group), or
  \item Use a window function (e.g., pick ''latest row'' with \texttt{ROW\_NUMBER}), or
  \item Aggregate intentionally (e.g., \texttt{MIN(created\_at)}), making the choice explicit.
\end{itemize}

\E PostgreSQL supports functional dependency optimizations in some cases (e.g., selecting 
columns functionally determined by a primary key), but you should not rely on this unless you 
fully understand the guarantees and your schema constraints.

\Pitfall
\begin{itemize}
  \item Using \texttt{DISTINCT} to ''fix'' grouping problems instead of correcting the grouping 
  logic or the joins that caused duplication.
\end{itemize}

\subsubsection{HAVING vs WHERE (filtering rows vs filtering groups)}
\J \texttt{WHERE} filters rows before grouping. \texttt{HAVING} filters groups after aggregation.

\M Logical order:
\begin{itemize}
  \item \texttt{WHERE} happens before \texttt{GROUP BY}.
  \item \texttt{HAVING} happens after aggregates are computed.
\end{itemize}
So \texttt{HAVING} can use aggregate expressions like \texttt{COUNT(*) > 10}.

\Sx Use \texttt{WHERE} for normal predicates whenever possible, because it reduces the amount 
of data that needs to be grouped. Reserve \texttt{HAVING} for conditions that truly depend on 
aggregates.

\E In performance tuning, a common win is moving predicates from \texttt{HAVING} to 
\texttt{WHERE} when they don't depend on aggregates.

\Pitfall
\begin{itemize}
  \item Putting simple row filters in \texttt{HAVING} can make queries significantly slower 
  and harder to understand.
\end{itemize}

\subsubsection{COUNT pitfalls (COUNT(*), COUNT(column), and NULLs)}
\J \texttt{COUNT(*)} counts rows. \texttt{COUNT(column)} counts non-NULL values in that column.

\M Because NULL represents ''unknown'', \texttt{COUNT(column)} ignores NULLs. This is often 
correct, but you must know what you're counting:
\begin{itemize}
  \item If you want ''number of rows'', use \texttt{COUNT(*)}.
  \item If you want ''number of rows where a field is present'', use \texttt{COUNT(column)}.
\end{itemize}

\Sx When counting distinct entities across joins, be explicit:
\begin{itemize}
  \item \texttt{COUNT(DISTINCT entity\_id)} avoids overcounting caused by join multiplicity.
  \item But \texttt{COUNT(DISTINCT ...)} can be expensive; fix join shape when possible.
\end{itemize}

\E For complex distinct counting, sometimes you restructure queries using CTEs/subqueries 
to reduce data volume before applying \texttt{COUNT(DISTINCT ...)}.

\Pitfall
\begin{itemize}
  \item Overcounting because joins multiply rows. Misinterpreting \texttt{COUNT(column)} 
  when NULLs exist.
\end{itemize}

\subsubsection{DISTINCT with aggregation (when it helps, when it hides bugs)}
\J \texttt{DISTINCT} removes duplicate rows from the result set.

\M In aggregation, you might use:
\begin{itemize}
  \item \texttt{COUNT(DISTINCT x)} to count unique values,
  \item or \texttt{SELECT DISTINCT ...} to remove duplicate rows.
\end{itemize}

\newpage
\Sx Treat \texttt{DISTINCT} as a tool, not a band-aid. If you need \texttt{DISTINCT} because 
the join or grouping logic creates unintended duplicates, fix the root cause first. Use 
\texttt{DISTINCT} when duplicates are semantically expected and you want a unique set.

\E In large datasets, \texttt{DISTINCT} can be expensive. For performance, you may prefer 
redesigning the query (pre-aggregation, dedup in a subquery, or correct join keys) rather than 
relying on \texttt{DISTINCT} late.

\Pitfall
\begin{itemize}
  \item Adding \texttt{DISTINCT} to ''make the numbers look right'' without understanding 
  where duplicates come from.
\end{itemize}

\subsubsection{Practical grouping patterns (the ones you actually use)}
\J Common patterns:
\begin{itemize}
  \item ''Total per entity'': \texttt{GROUP BY entity\_id}
  \item ''Totals by category'': \texttt{GROUP BY category}
  \item ''Totals over time'': \texttt{GROUP BY date\_trunc('month', created\_at)} (PostgreSQL)
\end{itemize}

\M Grouping over time often requires careful date handling (time zones, truncation) and 
clear definitions of the time bucket.

\Sx For time series in production, define a consistent time zone strategy and a clear 
''source of truth'' timestamp. This prevents subtle reporting mismatches.

\E For analytics-heavy needs, you may later add more advanced patterns (window functions, rollups) 
but the core is: define the grouping key precisely and keep it consistent.

\Pitfall
\begin{itemize}
  \item Grouping by a computed expression without understanding its stability (time zones, 
  rounding, formatting).
\end{itemize}